\section{Results and Discussion}

\subsection{Synthesis and Structural Characterization of Bi/MXene Composite}

The Bi/MXene composite was synthesized via a two-step route, as illustrated in Scheme~1. First, Ti$_{3}$C$_{2}$T$_{x}$ MXene nanosheets were prepared by selectively etching the Al layer from Ti$_{3}$AlC$_{2}$ MAX phase using a LiF/HCl etchant at 40~$^{\circ}$C for 48~h, followed by ice-bath ultrasonication to yield a high-concentration MXene colloidal suspension (20--25~mg~mL$^{-1}$). Subsequently, BiCl$_{3}\cdot x$H$_{2}$O (2~mmol) was dissolved in anhydrous ethanol and mixed with freeze-dried MXene aerogel (0.21~g) under magnetic stirring for 3~h. The mixture was transferred to a hydrothermal autoclave and heated at 180~$^{\circ}$C for 16~h, during which Bi$^{3+}$ hydrolyzed to form Bi(OH)$_{3}$ nanoparticles that nucleated on the MXene surface. The resulting Bi(OH)$_{3}$/MXene intermediate was collected by centrifugation, washed with ethanol and deionized water, and dried under vacuum at 60~$^{\circ}$C. Finally, the powder was calcined at 350~$^{\circ}$C under Ar atmosphere, during which Bi(OH)$_{3}$ underwent thermal decomposition and reduction to metallic Bi, yielding the Bi/MXene composite.

The morphology of the as-prepared materials was examined by scanning electron microscopy (SEM). Figure~Xa shows the pristine MXene nanosheets, which exhibit a characteristic layered morphology with pronounced surface wrinkles---features that provide an ideal platform for the uniform anchoring of Bi nanoparticles. In the Bi/MXene composite (Figure~Xb,c), abundant fine Bi nanoparticles are uniformly dispersed across the MXene surface and within the interlayer galleries, forming a distinctive sandwich-type architecture. This highly dispersed nanostructure is expected to shorten Na$^{+}$ diffusion pathways while the MXene layers mechanically encapsulate the Bi particles, buffering the severe volume fluctuations during sodiation/desodiation and thereby enhancing electrode integrity.

Energy-dispersive X-ray spectroscopy (EDS) elemental mapping (Figure~X) provides direct evidence for the confined distribution of Bi within the MXene scaffold. The characteristic MXene elements---Ti, O, and F---exhibit strong spatial correlation with the Bi signal, confirming that Bi nanoparticles are intimately integrated into the MXene layered matrix. This tight interfacial contact constitutes the physical foundation of the sandwich confinement architecture, enabling MXene to impose robust mechanical constraint on Bi during cycling and thereby substantially improving structural stability and cycle life.

The atomic-scale interfacial structure of the Bi/MXene composite was further elucidated by high-resolution transmission electron microscopy (HRTEM). As shown in Figure~X, Bi nanocrystals are observed to be tightly encapsulated by MXene nanosheets. The well-resolved lattice fringes with an interplanar spacing of 0.329~nm correspond to the (012) plane of trigonal Bi, consistent with the XRD results. This intimate Bi--MXene interfacial contact provides direct morphological evidence for the sandwich confinement architecture. In this configuration, the MXene layers serve a dual structural function: they act as a highly conductive substrate that accelerates electron transport, while their mechanically robust two-dimensional framework exerts strong physical confinement on the Bi particles. This nanoscale confinement effect, realized through the intimate Bi--MXene interfacial contact, fundamentally suppresses particle pulverization during the sodiation/desodiation process, thereby imparting exceptional structural stability and long-term cycling durability to the electrode.

The crystal structure of the Bi/MXene composite was characterized by X-ray diffraction (XRD). As shown in Figure~2a, a strong diffraction peak appears at $2\theta = 27.2^{\circ}$, which matches the (012) plane of trigonal-phase Bi (JCPDS no.\ 85-1329), confirming that Bi is present in a well-crystallized metallic state. Concurrently, a characteristic (002) reflection of MXene is observed at $2\theta \approx 8^{\circ}$, indicating that the layered structure of Ti$_{3}$C$_{2}$T$_{x}$ is preserved after the hydrothermal and calcination treatments.

To further probe the interfacial chemical interaction between Bi and MXene, X-ray photoelectron spectroscopy (XPS) was performed. The survey spectra (Figure~3a) confirm the coexistence of Ti, O, F, C, and Bi in both pristine MXene and the Bi/MXene composite. In the high-resolution Bi~4f spectrum of Bi/MXene (Figure~3d), two prominent symmetric peaks located at 164.56 and 159.26~eV correspond to metallic Bi (Bi~4f$_{5/2}$ and Bi~4f$_{7/2}$), while weaker signals at 163.56 and 157.76~eV indicate the presence of a thin native oxide layer on the Bi nanoparticle surface, providing the requisite active sites for subsequent chemical bonding. The most critical evidence emerges from comparative analysis of the O~1s and Ti~2p spectra (Figure~3b,c vs.\ 3e,f). In pristine MXene, the O~1s spectrum is dominated by Ti--O (530.14~eV), Ti--OH (530.94~eV), and adsorbed water. In the Bi/MXene composite, however, a prominent new component appears at 527.3~eV in the O~1s spectrum, assignable to Ti--O--Bi covalent bonds. Concomitantly, the Ti~2p spectrum (Figure~3f) exhibits a slight binding energy shift of the Ti--O-related peaks and a new component at 456.5~eV corresponding to Ti--O--Bi linkages. The synchronous emergence of this signature in two independent core-level spectra strongly indicates that Bi nanoparticles are covalently anchored to the MXene surface through reaction with oxygen-containing terminations, rather than merely physically adsorbed. This pronounced interfacial chemical interaction significantly enhances the Bi--MXene adhesion, providing a chemical safeguard that complements the physical confinement, thereby fundamentally suppressing particle pulverization during cycling and imparting exceptional long-term structural stability.

\subsection{Electrochemical Sodium Storage Performance}

The electrochemical sodium storage behavior of the Bi/MXene composite was evaluated in half-cell configurations. The rate capability was assessed by measuring the specific capacity at current densities ranging from 0.1 to 20~A~g$^{-1}$ (Figure~4a). The electrode delivers reversible capacities of 356.6, 354.8, 351.6, 348.9, 346.2, 341.8, 336.9, and 316.9~mAh~g$^{-1}$ at 0.1, 0.2, 0.5, 1, 2, 5, 10, and 20~A~g$^{-1}$, respectively. Notably, increasing the current density by a factor of 200 (from 0.1 to 20~A~g$^{-1}$) results in only an 11.1\% reduction in capacity, underscoring the exceptional rate tolerance of the Bi/MXene composite. Furthermore, when the current density is returned to 0.1~A~g$^{-1}$, the specific capacity recovers to 362.9~mAh~g$^{-1}$ — slightly exceeding the initial value of 356.6~mAh~g$^{-1}$, indicative of progressive electrode activation during the rate test. This complete capacity recovery after exposure to high-rate cycling demonstrates the outstanding structural robustness of the Bi/MXene electrode. In contrast, both pristine Bi and MXene electrodes exhibit poor rate capability, with negligible capacity contribution at high current densities (Figure~S). This exceptional rate performance is attributed to the rapid electron transport network provided by the MXene scaffold and the shortened Na$^{+}$ diffusion distance afforded by the nanoscale Bi particles.

Figure~4b presents the galvanostatic charge--discharge (GCD) voltage profiles at selected cycle numbers. During the discharge process, a well-defined voltage plateau is observed, corresponding to the alloying reaction between Bi and Na to form Na$_{3}$Bi. Notably, the charge--discharge curves from the 3rd to the 100th cycle are nearly superimposable, indicating exceptionally high electrochemical reversibility of the Bi/MXene electrode. Furthermore, the remarkably small voltage hysteresis between the charge and discharge branches provides direct evidence for the rapid ion transport kinetics and low charge-transfer resistance at the Bi/MXene interface. Cycling performance at 1~A~g$^{-1}$ (Figure~4c) reveals a reversible capacity of 333.8~mAh~g$^{-1}$ after 100 cycles, corresponding to 98.5\% capacity retention (relative to the 3rd cycle).

The long-term cycling stability of the Bi/MXene composite was assessed at a current density of 2~A~g$^{-1}$ (Figure~4d). In the first cycle, the electrode delivers a discharge capacity of 422.9~mAh~g$^{-1}$ and a charge capacity of 372.6~mAh~g$^{-1}$, yielding an initial Coulombic efficiency (ICE) of 88.1\%. The irreversible capacity loss of approximately 50~mAh~g$^{-1}$ in the first cycle is attributable to the formation of the solid-electrolyte interphase (SEI) and the trapping of Na$^{+}$ in residual surface functional groups of MXene. After an initial activation period of three cycles, the electrode stabilizes at a reversible capacity of 343.0~mAh~g$^{-1}$ (4th cycle). Upon extended cycling, the electrode exhibits outstanding durability, retaining a specific capacity of approximately 310~mAh~g$^{-1}$ after 5,000 consecutive charge--discharge cycles, corresponding to a capacity retention of 90.4\% relative to the stabilized 4th-cycle capacity. The Coulombic efficiency rises rapidly after the initial cycles and stabilizes near 100\% throughout the remaining cycling period. This remarkable cycling stability is ascribed to the sandwich-type confinement effect of the MXene scaffold, which effectively accommodates the severe volume changes of Bi during sodiation/desodiation, suppresses active material pulverization and delamination, and preserves the long-term integrity of the electrode architecture. Notably, the rate capability of the Bi/MXene anode compares favorably with other reported Bi-based anodes (Figure~4e) [10,11,17].

\subsection{Reaction Kinetics and Charge Storage Mechanism}

To elucidate the charge storage mechanism underlying the exceptional rate performance, cyclic voltammetry (CV) was performed at scan rates ranging from 0.2 to 5.0~mV~s$^{-1}$ (Figure~X). The well-defined redox peaks at all scan rates confirm the highly reversible alloying/dealloying reaction of Bi with Na. The relationship between peak current ($i$) and scan rate ($v$) was analyzed according to the power-law equation $i = av^{b}$, where the $b$-value distinguishes between diffusion-controlled Faradaic processes ($b \approx 0.5$) and surface-confined capacitive processes ($b \approx 1.0$). Linear fitting of $\log i$ versus $\log v$ yields $b$-values in the range of 0.62--0.82, indicating that the charge storage proceeds through a hybrid mechanism involving both diffusion-controlled battery-type behavior and surface-confined pseudocapacitive contributions. Quantitative deconvolution of the current response reveals that the diffusion-controlled process dominates, accounting for 92\%--98\% of the total capacity across the measured scan rates. Nevertheless, the non-negligible capacitive contribution, facilitated by the MXene scaffold, plays a critical role in sustaining rapid charge transfer at high current densities, thereby enabling the electrode to maintain a high capacity under extreme rate conditions.

Electrochemical impedance spectroscopy (EIS) was conducted to probe the charge transport properties of the Bi/MXene composite in comparison with a pristine Bi electrode (Figure~X). The Nyquist plot of the Bi/MXene electrode exhibits a dramatically smaller semicircle in the high-to-medium frequency region relative to that of pure Bi, reflecting a substantially reduced charge-transfer resistance ($R_{\text{ct}}$ = 0.69~$\Omega$) relative to that of pure Bi ($R_{\text{ct}}$ = 13.95~$\Omega$). In addition, the Bi/MXene electrode displays a steeper Warburg slope in the low-frequency region, indicative of more facile Na$^{+}$ diffusion within the electrode. These results provide compelling evidence that the two-dimensional MXene conductive framework, in intimate contact with Bi nanoparticles, establishes highly efficient electron transport pathways that significantly lower the interfacial charge-transfer barrier and accelerate the electrode reaction kinetics.

The galvanostatic intermittent titration technique (GITT) was employed to quantitatively evaluate the Na$^{+}$ diffusion coefficient ($D_{\text{Na}^{+}}$) across the full sodiation/desodiation voltage range (Figure~X). The Bi/MXene electrode exhibits a high $D_{\text{Na}^{+}}$ throughout the entire voltage window, with values in the range of $10^{-12.5}$--$10^{-13}$~cm$^{2}$~s$^{-1}$. Notably, the polarization voltage during the GITT measurement is as low as 74.6~mV, reflecting minimal ohmic and concentration overpotentials. This result directly validates the sandwich confinement design principle: the MXene scaffold enhances electronic conductivity, while the nanoscale Bi particles drastically shorten the solid-state Na$^{+}$ diffusion length, enabling a synergistic rapid ion/electron co-transport at the atomic scale.

\subsection{Post-Cycling Structural Analysis}

To elucidate the intrinsic mechanism underlying the exceptional cycling stability of the Bi/MXene composite, we constructed a mechanistic model as depicted in Figure~6a. For a pristine Bi anode, Bi particles undergo severe volume expansion and contraction during consecutive sodiation/desodiation cycles (upper panel of Figure~6a). The accumulated mechanical stress induces irreversible pulverization and structural collapse, leading to electrical contact failure and rapid capacity decay. In stark contrast, Bi nanoparticles anchored within the MXene interlayer galleries form a distinctive sandwich confinement architecture. In this configuration, MXene plays a dual role: the mechanically robust MXene nanosheets provide strong physical confinement that effectively buffers and suppresses the volume expansion of Bi particles, preventing particle fracture (lower panel and inset of Figure~6a); meanwhile, the high electrical conductivity of MXene offers rapid electron transport pathways, and the shortened Na$^{+}$ diffusion distance within the nanoscale-confined architecture significantly enhances electrochemical reaction kinetics. This synergistic ``physical confinement--kinetic acceleration'' mechanism fundamentally ensures the structural integrity of the electrode over ultra-long cycling.

To assess the structural integrity of the Bi/MXene electrode after prolonged electrochemical cycling, ex-situ SEM was performed on the electrode after 5,000 cycles at 2~A~g$^{-1}$ (Figure~6b,c). In stark contrast to the severe pulverization and agglomeration typically observed in cycled Bi anodes [9,10], the Bi/MXene electrode largely retains its original morphology, with Bi nanoparticles remaining well-dispersed within the MXene interlayer galleries and no evidence of extensive cracking or delamination. This morphological stability provides direct visual confirmation that the sandwich-type MXene confinement effectively suppresses the mechanical degradation pathways — particle pulverization, electrical disconnection, and SEI re-formation — that conventionally plague alloy-type anodes. These post-cycling observations, together with the kinetic analyses, collectively establish that the Bi/MXene sandwich architecture enables robust electrode-level structural stability while maintaining rapid ion/electron transport, thereby underpinning the simultaneous achievement of high rate capability and long cycling durability.

\section{Conclusion}

In summary, we have demonstrated a binary Bi/MXene composite anode with a distinctive sandwich architecture, synthesized via a facile hydrothermal route followed by thermal reduction. The Bi nanoparticles are uniformly confined between Ti$_{3}$C$_{2}$T$_{x}$ MXene nanosheets, as confirmed by SEM, HRTEM, EDS elemental mapping, XRD, and XPS analyses. XPS further reveals the formation of covalent Ti--O--Bi bonds at the Bi--MXene interface, demonstrating that the Bi nanoparticles are chemically anchored to the MXene surface, providing a chemical safeguard that reinforces the physical confinement. This structurally simple design, requiring no sulfidation, carbon coating, or multi-component heterostructure engineering, simultaneously delivers exceptional rate capability and ultra-long cycling stability.

The Bi/MXene electrode achieves a reversible capacity of 356.6~mAh~g$^{-1}$ at 0.1~A~g$^{-1}$ and retains 316.9~mAh~g$^{-1}$ at 20~A~g$^{-1}$, corresponding to 88.8\% capacity retention over a 200-fold increase in current density. Upon extended cycling, the electrode maintains approximately 310~mAh~g$^{-1}$ after 5,000 cycles at 2~A~g$^{-1}$ with a capacity retention of 90.4\% and a Coulombic efficiency approaching 100\%. The initial Coulombic efficiency reaches 88.1\%.

Mechanistic investigations reveal that the superior electrochemical performance arises from a synergistic combination of physical confinement and chemical anchoring. The MXene scaffold provides a continuous metallic conductive network that reduces the charge-transfer resistance from 13.95~$\Omega$ (pure Bi) to 0.69~$\Omega$, while its mechanically robust two-dimensional framework physically confines the Bi nanoparticles. The covalent Ti--O--Bi bonds chemically anchor Bi to the MXene surface, collectively suppressing the pulverization and electrical disconnection that conventionally cause capacity fading in alloy-type anodes. CV kinetic analysis confirms a hybrid charge storage mechanism ($b$ = 0.62--0.82), in which the pseudocapacitive contribution from MXene sustains rapid charge transfer at high rates, complementing the dominant diffusion-controlled capacity of Bi. GITT measurements corroborate these findings, yielding Na$^{+}$ diffusion coefficients of $10^{-12.5}$--$10^{-13}$~cm$^{2}$~s$^{-1}$ with a remarkably low polarization voltage of 74.6~mV. Post-cycling SEM further verifies that the sandwich architecture preserves electrode morphology after 5,000 cycles, consistent with the confinement-driven structural stability central to our design.

These results establish that the rate-capacity-stability trade-off long recognized in alloy-type SIB anodes can be overcome through a structurally simple, scalable binary composite architecture — opening a practical pathway toward high-performance fast-charging sodium-ion batteries.
